\subsection{The characteristic-two K3 possibility}

We now describe the K3 case permitted for an ample divisor. The splitting
below follows from duality on the singular surface itself. It does
not require a positive-characteristic vanishing theorem.

\begin{lemma}\label{lem:k3-ample-curve-vanishing}
Let $V$ be a projective integral Gorenstein surface over an
algebraically closed field, with $\omega_V\simeq\OO_V$ and
$H^1(V,\OO_V)=0$. If $C\subset V$ is a smooth ample integral
Cartier curve, then
\[
 H^1(V,\OO_V(nC))=0\quad(n\in\Z\setminus\{0\}),
 \qquad h^0(V,\OO_V(C))=1+g(C)=2+C^2/2.
\]
The line bundle $\OO_V(C)$ is globally generated.
\end{lemma}

\begin{proof}
Adjunction identifies $\OO_C(C)$ with $\omega_C$. The sequence
\[
 0\longrightarrow\OO_V\longrightarrow\OO_V(C)
 \longrightarrow\omega_C\longrightarrow0
\]
therefore gives the asserted dimension of sections. By Serre
duality, $H^2(\OO_V(C))=H^0(\OO_V(-C))^*=0$, whereas
$H^2(\OO_V)$ and $H^1(\omega_C)$ both have dimension one.
The surjection $H^1(\omega_C)\to H^2(\OO_V)$ is consequently
an isomorphism. Its kernel is $H^1(\OO_V(C))$, which vanishes.

For $n\geq2$, use the divisor sequence with successive twists
$(n-1)C$ and $nC$. Curve duality gives
\[
 H^1(C,\OO_C(nC))=H^0(C,\OO_C((1-n)C))^*=0.
\]
Induction proves the positive twists; surface duality proves the
negative ones. These applications of adjunction and duality are
valid for a projective Gorenstein surface
\cite[Tags 0AA4, 0FVZ and 0FW0]{Stacks}.

Finally $C^2=2g(C)-2>0$, so $g(C)\geq2$. The canonical system
of $C$ has no base point. Indeed, for a point $x\in C$ one has
$h^0(\OO_C(x))=1$, since a degree-one map to $\PP^1$ would
make $C$ rational; Riemann--Roch gives
$h^0(\omega_C(-x))=g(C)-1$. The surjection of sections onto
$H^0(\omega_C)$ proves generation along $C$, and the defining
section of $C$ proves it away from $C$.
\end{proof}

\begin{proposition}\label{prop:k3-picard-derivations}
In the non-almost-ruled alternative of
Proposition~\ref{prop:normalized-mup-ekedahl}, put $s=Y^2$.
The tangent sequence has a unique splitting, and
\[
 T_N\simeq\OO_N(-D)\oplus\OO_N(D),\qquad
 h^0(N,T_N)=2+s,\qquad H^1(N,T_N)=0.
\]
All global derivations lie in the positive summand and generate
it. Every nonzero global derivation $v$ satisfies $v^2=c_vv$ for
some $c_v\in k$. Its normal quotient is rational. There are
dominant purely inseparable rational maps from rational surfaces
to $N$ of degree two and to $X$ of degree four.
\end{proposition}

\begin{proof}
Rational double points are rational and crepant singularities.
The K3 resolution therefore gives $\omega_N\simeq\OO_N$ and
$H^1(\OO_N)=0$. The extension class of
\eqref{eq:normalized-mup-tangent-sequence} belongs to
\[
 \operatorname{Ext}^1(\OO_N(D),\OO_N(-D))
 =H^1(\OO_N(-2D))=0
\]
by Lemma~\ref{lem:k3-ample-curve-vanishing}. The splitting is
unique because $H^0(\OO_N(-2D))=0$. The lemma also gives
the cohomology dimensions and global generation of the positive
summand $P=\OO_N(D)$, since $g(D)=s+1$.

For a local generator $e$ of $P$, the class of $e^2$ modulo $P$
transforms quadratically: in characteristic two,
$(fe)^2=f^2e^2+f e(f)e$. Thus it defines a homomorphism
\[
 P^{\otimes2}\longrightarrow T_N/P=\OO_N(-D).
\]
It is zero because $H^0(\OO_N(-3D))=0$. Write $e^2=he$
with $h$ regular. If a global derivation is locally $v=fe$, then
\[
 v^2=(fh+e(f))v.
\]
The coefficient is regular even at the zeros of $f$ and agrees
in different trivializations on the dense locus where $v\ne0$.
It therefore defines a global regular function on $N$, hence a
scalar $c_v$. If $c_v\ne0$, rescale $v$ so that $v^2=v$;
otherwise $v^2=0$. These are respectively a $\mu_2$-action and
an $\alpha_2$-action
\cite[Proposition~2.10, author version p.~5]{MatsumotoActions}.

At every codimension-one point, $N$ is smooth and $P$ is a line
subbundle of its tangent bundle. The divisorial fixed locus of
$v$ is therefore the zero divisor of the corresponding section
of $P$. It is linearly equivalent to $D$, and is nonzero.
The quotient $R=N^v$ is normal and rational by
\cite[Proposition~4.1 and proof, author version pp.~16--17]{MatsumotoActions}:
the nonzero divisorial fixed part places it in the rational
alternative there.

For the degree assertion write $E=k(N)$ and $J=k(R)$.
The quotient has $[E:J]=2$, while $E^2\subset J$ and
$[E:E^2]=4$. Hence $E\subset J^{1/2}$ and
$[J^{1/2}:E]=2$. The field $J^{1/2}$ is rational over the
perfect ground field. It gives a degree-two rational domination
of $N$, and composition with $N\to X$ gives degree four.
\end{proof}

The degree-four conclusion does not imply that $X$ is almost
ruled: that term here requires degree two. In particular the
large space of global derivations does not exclude this case.

\begin{proposition}\label{prop:k3-picard-paracanonical}
In characteristic two, in any equality case
$2g(Y)-2=2aY^2$ of
Proposition~\ref{prop:normalized-mup-conductor}, one has
\[
 \pi_*\OO_N/\OO_X\simeq M^{-1}(aY).
\]
In the non-almost-ruled alternative of
Proposition~\ref{prop:normalized-mup-ekedahl}, the cover
$N\to X$ is an $\alpha_{\omega_X^{-1}}$-torsor and
\[
 0\longrightarrow\OO_X\longrightarrow\pi_*\OO_N
 \longrightarrow\omega_X\longrightarrow0,
 \qquad M\simeq\OO_X(Y)\otimes\omega_X^{-1}.
\]
Thus $Y$ belongs to $|\omega_X\otimes M|$, with $M$ a
nontrivial two-torsion twist, and
$\OO_N(D)\simeq\pi^*\omega_X$.
\end{proposition}

\begin{proof}
The unit in the locally free rank-two algebra $\pi_*\OO_N$
is nonzero on every fibre. Its quotient $Q$ is therefore an
invertible sheaf. The inclusion of the Kummer algebra
$\OO_X\oplus M^{-1}$ gives a nonzero map $M^{-1}\to Q$,
an isomorphism away from $Y$. At the generic point of $Y$,
the normalized basis in the conductor proof is $1,w$, with
$w=(z-v)/t^a$. Modulo the base summand, $[z]=t^a[w]$.
The divisor of the map of line bundles is exactly $aY$;
there is no additional isolated-zero divisor. Hence
$Q=M^{-1}(aY)$.

Now assume the non-almost-ruled case. Since $K_X$ is ample,
$X$ is its own minimal model. Ekedahl's theorem gives a normal
$\alpha_{\omega_X^{-1}}$-torsor $C\to X$ with rational double
points and K3 resolution
\cite[I, Theorem~2.4(ii)(d), p.~109]{EkedahlCanonical}.
We identify this cover with $N$. For an $\alpha_L$-torsor,
the differential of its equation gives $L^{-2}\to\Omega^1_X$.
Normality of $C$ makes its zero scheme finite, so its saturated
annihilator is $L^{-2}\otimes\omega_X^{-1}\subset T_X$.
For $L=\omega_X^{-1}$ this is an embedding $\omega_X\to T_X$.
The foliation belonging to $N$ has the same source because
$\OO_X(2Y)=\omega_X^2$.

Any two such embeddings have zero wedge: the wedge is a section
of $\operatorname{Hom}(\omega_X^2,\det T_X)=\omega_X^{-3}$,
which has no nonzero section. The two generic lines agree and,
by saturation, define the same foliation. Its invariant field
determines the square of the covering field. Uniqueness of square
roots in characteristic two therefore identifies the two field
extensions over $k(X)$, and normality identifies $C$ and $N$
over $X$, as in the proof of
Proposition~\ref{prop:normalized-mup-ekedahl}.

The algebra filtration for this additive torsor has quotient
$\omega_X$ \cite[I, Proposition~1.7, p.~107]{EkedahlCanonical}.
Comparison with $Q=M^{-1}(Y)$ gives the asserted formula for
$M$. The pullback of $M$ to its Kummer torsor is trivial;
pulling this trivialization back along $N\to T$ gives
$\pi^*M\simeq\OO_N$. The final identity follows.
\end{proof}

\begin{proposition}[After Matsumoto]\label{prop:k3-picard-singularities}
In the non-almost-ruled alternative of
Proposition~\ref{prop:normalized-mup-ekedahl}, every singularity
of $N$ has type
\[
 A_1,\qquad D_{2m}^{0}\ (m\geq2),\qquad E_7^{0},\qquad E_8^{0}.
\]
The sum of the Dynkin indices is $12+Y^2$, between thirteen
and twenty-one. This is a necessary condition; it does not assert
that every configuration with this sum occurs.
\end{proposition}

\begin{proof}
All singularities lie away from $D$, where $N=T$ is the original
$\mu_2$-torsor over smooth $X$. Its derivation has no fixed
point there. The smooth-quotient rows in
\cite[Theorem~3.3 and Table~2, author version pp.~10--11]{MatsumotoActions}
give exactly the listed types. The local multiplicative
classification is proved in
\cite[Proposition~4.8 and Table~3, pp.~13--15]{MatsumotoMuN}.
The superscript zero specifies the characteristic-two local type.

The index sum can also be read directly from their equations.
Write the local equation as $z^2+h(x,y)=0$. The Fitting
singular ideal is $(h_x,h_y)$ in a free rank-two algebra over
$k[[x,y]]$. The relevant base lengths are
\[
\begin{array}{c|c|c}
 \text{type}&h&\dim_k k[[x,y]]/(h_x,h_y)\\ \hline
 A_1&xy&1\\
 D_{2m}^{0}&x^2y+xy^m&2m\\
 E_7^{0}&x^3+xy^3&7\\
 E_8^{0}&x^3+y^5&8.
\end{array}
\]
For the $D$ row, the quotient is by
$(y^m,x^2+\epsilon xy^{m-1})$, where $\epsilon=m\bmod2$;
it is free of rank two over $k[[y]]/(y^m)$. For the $E_7$
row, a basis is $1,y,y^2,y^3,y^4,x,xy$. For the $E_8$
row, the ideal is $(x^2,y^4)$. Each Fitting length upstairs is
twice the indicated Dynkin index. Their sum is
$2(12+Y^2)$ by
Proposition~\ref{prop:normalized-mup-ekedahl}, which proves
the assertion.
\end{proof}

For a merely ample divisor these constraints leave a specific
existence question: can such
a characteristic-two surface have a smooth paracanonical curve
on which its nontrivial two-torsion twist becomes trivial? The
splitting, rational domination and local singularity list do not
answer that question.


Gounelas and Liedtke prove that a smooth K3 surface has no nonzero
global symmetric differential of positive degree except in
characteristic two with Artin invariant one
\cite[Theorem~1.1, p.~2]{GounelasLiedtkeSymmetric}.
In the exceptional case the dimension is one in positive even
degrees and zero in odd degrees. Their theorem concerns the
ordinary cotangent bundle of the smooth resolution. Here the
splitting gives the reflexive cotangent bundle on the singular
surface $N$. Applying that theorem would require proving that a
specified nonzero symmetric tensor pulls back without poles along
every exceptional component. The splitting and the singularity
list alone do not provide this extension, so they do not reduce
the remaining alternative to Artin invariant one.
